\ssscp{\tsc{West Flores}}{03-02-03-01}
Within \tsc{West Central Flores},
Manggarai, Rembong, Waerana, and Komodo can be placed together
in a single node
and apart from \tsc{Central Flores} on the basis of shared *z
> \it{d} in at least four words,
further merger of *d with *z/*j > **z in six words,
and shared *-aw > \it{o} in most cases.
While we currently treat Komodo as a member of \tsc{West Flores},
recall from \srf{03-01-01} that Komodo also shares some sound changes with Bima.
This may indicate that Bima
and Komodo belong in a single node,
and/or that the features shared between Komodo
and \tsc{West Flores} are due to contact.

Examples of *z > \it{d} in West Flores are given in \trf{03-30}.
Where reflexes are available,
most \tsc{Central Flores} languages have the same reflexes as
attested for the merger of *z/*j discussed above.
However, Palu{\Q}e has one case of *z > \it{t} (*zauq > \it{teu}
‘far’) and Lio *zaqat > \it{reʔe} has *z > \it{r}.

\begin{table}\tcp{\tsc{West Flores} *z > \it{d}}{03-29}
	\begin{threeparttable}\st{6pt}
	\begin{tabular}{lllll}\lsptoprule
*gloss	&	far	&	evil	&	feel, grope	&	ladder\\
PMP	&	*zauq	&	*zaqat	&	*zəmak	&	*haRəzan\\	\midrule
Komodo	&	‹\tbr{d}iu›\su{†}	&	‹\tbr{d}aaʔ›	&	‹\tbr{d}ama›	&	\cg{(raŋge)}\\
Manggarai	&	{\hr}\tbr{d}eu	&	{\hr}\tbr{d}aʔat	&	{\hr}\tbr{d}əmak	&	{\hr}rə\tbr{d}aŋ, rə\tbr{d}e\\
Kepo{\Q}	&	{\hr}\tbr{d}euʔ	&	{\hr}\tbr{d}aʔat	&		&	\\
Waerana	&	{\hr}\tbr{d}euʔ	&	{\hr}\tbr{d}aʔat	&	{\hr}\tbr{d}əma \it{‘taste’}	&	\cg{(radi)}\su{‡}\\
Rembong	&	{\hr}\tbr{d}iuʔ	&	{\hr}\tbr{d}aʔat	&	{\hr}\tbr{d}amak	&	{\hr}rə\tbr{d}e\\	\midrule
Palu{\Q}e	&	{\hr}teu	&	\cg{(ndoa)}	&		&	\\
Lio	&	\cg{(beu)}	&	{\hr}reʔe	&		&	\cg{(taŋi)}\\
Ende	&	{\hr}reu	&	{\hr}reʔe	&	\cg{(ɗəpa)}	&	\cg{(ŋgera)}\\
C. Nage	&	{\hr}zeu	&	{\hr}eʔe	&		&	\\
Keo, Ud.	&	{\hr}reu	&	{\hr}reʔe	&	\cg{(dede)}	&	\cg{(taŋi)}\\
Ngadha	&	{\hr}ʣeu	&	{\hr}ʣeʔe	&		&	\cg{(raɗi, taŋi)}\\
Rongga	&	{\hr}ɹeu	&	{\hr}ɹeʔe	&		&	\cg{(raɗi)}\\
		\lspbottomrule
	\end{tabular}
		\begin{tablenotes}\tnsize
			\item[†]	{Recall from \srf{03-02-01} that Komodo has a distinction between
								/d/ and /ɗ/ which is not represented in \citet{Verheijen1982}.
								Komodo words with ‹d› are enclosed in angles to indicate this under-differentiation.}
			\item[‡]	{Evidence that Waerana \it{radi} is not a reflex of *haRəzan
								comes from the alternate Rembong word \it{radi} which stands
								alongside \it{rəde} from *haRəzan.}			
		\end{tablenotes}
	\end{threeparttable}
\end{table}

Examples of the partial merger of *d with *z/*j in \tsc{West
Flores} are given in \trf{03-30}.
While *d merges with *z/*j in most \tsc{Central Flores} languages,
it is distinct in Lio which has *d > \it{r} but *z/*j > \it{ʤ}.
While contact may have played a role in the spread of *d/**z
> **z in \tsc{West Flores}
and \tsc{Central Flores}, we can posit that such contact took
place before the break up of Proto-West Flores.
Thus, partial merger of *d with *d/*j is supporting evidence
for \tsc{West Flores}.

\newpage
\begin{table}\tcp{\tsc{West Flores} partial merger of *d/*z/*j}{03-30}
	\begin{threeparttable}\st{4pt}
	\begin{tabular}{llllll@{}l}\lsptoprule
*gloss	&	leaf	&	two	&	stand	&	behind	&	young	&	knee\\
PMP	&	*\tbr{d}ahun	&	*\tbr{d}uha	&	*kə\tbr{d}əŋ	&	*ma-u\tbr{d}əhi	&	*ŋu\tbr{d}a	&	*tuqu\tbr{d}\\	\midrule
Komodo	&	\cg{(ŋolo)}	&	{\hr}\tbr{r}ua	&	{\hr}hə\tbr{r}e	&	\cg{(kuni)}	&	{\hr}ŋu\tbr{r}a	&	{\hr}tuu\tbr{h}\su{‖}\\
Manggarai	&	{\hr}\tbr{s}auŋ	&	{\hr}\tbr{s}ua	&	{\hr}hə\tbr{s}e	&	{\hr}mu\tbr{s}i	&	{\hr}ŋu\tbr{s}e\su{‡}	&	{\hr}tuʔu\tbr{s}\\
Kepo{\Q}	&	{\hr}\tbr{r}aun	&	{\hr}\tbr{r}ueʔ	&	\cg{(ndaiŋ)}	&	{\hr}mu\tbr{r}i	&		&	{\hr}tuʔu\tbr{s}\\
Waerana	&	{\hr}\tbr{z}aun\su{†}	&	{\hr}\tbr{z}uaʔ	&	\cg{(nadir)}	&	{\hr}mu\tbr{z}i	&	{\hr}ŋu\tbr{z}a	&	\cg{(lulur)}\\
Rembong	&	{\hr}\tbr{z}aun\su{†}	&	{\hr}\tbr{z}ua	&	{\hr}gə\tbr{z}eʔ	&	{\hr}mu\tbr{z}i	&	{\hr}ŋu\tbr{z}a	&	{\hr}tuʔu\tbr{s}\\	\midrule
Namut	&	\cg{(wunu)}	&	{\hr}rua	&	\cg{(ndai)}	&		&		&	{\hr}ulu tuu\\
Palu{\Q}e	&	\cg{(vunu)}	&	{\hr}rua	&	{\hr}kəre	&	{\hr}muri	&		&	{\hr}tuu\\
Lio	&	\cg{(ʋunu)}	&	{\hr}rua	&	\cg{(dari)}	&	{\hr}muri	&	{\hr}ŋura, ŋuʤa	&	{\hr}mbuku tuu\\
Ende	&	\cg{(wunu)}	&	{\hr}rua	&	\cg{(dari)}	&	{\hr}muri	&	{\hr}ŋura	&	\cg{(nduku)}\\
C. Nage	&	\cg{(wunu)}	&	{\hr}ɗua	&	\cg{(dawi)}	&	{\hr}muzi	&		&	{\hr}ulu tuu\\
Keo, Ud.	&	\cg{(wunu)}	&	{\hr}rua	&	\cg{(ndeli)}	&	{\hr}muri	&	{\hr}ŋura	&	{\hr}ʔudu tuu\\
Ngadha	&	\cg{(vunu)}	&	{\hr}ʣua	&	\cg{(duge)}	&	{\hr}muʣi	&	{\hr}ŋura, ŋuʣa\su{Δ}	&	{\hr}ulu tuu\\
Rongga	&	\cg{(wunu)}	&	{\hr}ɹua	&	\cg{(ndəri)}	&	{\hr}muɹi	&	{\hr}ŋuɹa	&	{\hr}ulu tuu\\
		\lspbottomrule
	\end{tabular}
		\begin{tablenotes}\tnsize
			\item[†]	{Rembong \it{zaun} = ‘leaf husk (e.g. corn to smoke)’.
								Waerana \it{zaun} = ‘husk (corn), leaf’.}
			\item[‡]	{\it{ŋuse} ‘young’ is from Riung in the etymological notes in \citet[114]{Verheijen1982}.
								Riung is classified by \citet{Verheijen1967} as a Manggarai dialect.}
			\item[Δ]	{Ngadha \it{ŋura,
								ŋuʣa} is from \citet[114]{Verheijen1982}.}
			\item[‖]	{Final \it{h} in Komodo \it{tuuh} ‘knee’ could be from intermediate **s or **r.
								Several words are listed in \citet{Verheijen1982} which have
								alternates ending in \it{s {\tl} h} as well as \it{r {\tl} h}.}
		\end{tablenotes}
	\end{threeparttable}
\end{table}

We cannot posit complete merger of *d/*z/*j in \tsc{West Flores},
as some words require **d to be distinct from *z/*j at Proto-West Flores.
Examples of *d = **d are given in \trf{03-31},
in which reflexes that are the same as *z/*j are shaded.

Another change shared between \tsc{West Flores} languages which
sets them apart from \tsc{Central Flores} is final
*-aw > \it{o} in most words.
\tsc{Central Flores}, on the other hand, usually has final *-aw > \it{a}.
Final *-aw > \it{o} also occurs in \tsc{East Sumbawa}
and \tsc{Havu-Dhao} (\srf{03-01}) and thus this is not an
exclusively shared sound change.
Examples of PMP *-aw in \tsc{West Flores} are given in \trf{03-32}.

The changes affecting final *-ay in West Flores languages are more diverse.
Komodo has *-ay > \it{e},
Rembong and Manggarai usually have *-ay > \it{a},
and Waerana has *-ay > \it{ar}.
Waerana *-ay > \it{ar} may have gone through intermediate **-aʤ
and/or **-az, though complicating this hypothesis are the two cases of final
*{\nbhyp}aw > \it{ar} in \trf{03-32}.
Other reflexes of final *-ay are also occasionally attested.
Examples of final *-ay in the \tsc{West Flores} languages are
given in \trf{03-33}.

\begin{table}[p]\centering\tcp{\tsc{West Flores} *d = **d}{03-31}
\begin{adjustbox}{width=\textwidth}
	\begin{threeparttable}\st{2pt}
	\begin{tabular}{llllll|lll}\lsptoprule
*gloss	&	blood	&	good	&	deep	&	bathe	&	house post	&	body dirt	&	live\\
PMP	&	*\tbr{d}aRaq	&	*\tbr{d}iqa	&	*\tbr{d}aləm	&	*\tbr{d}iRus	&	*ha\tbr{d}iRi	&	*\tbr{d}aki	&	*ma-qu\tbr{d}ip\\ \midrule
Komodo	&	{\hr}\tbr{r}aa	{\cgr}&	\cg{(reheŋ)}	&	{\hr}\tbr{d}oleŋ	&	\cg{(buri)}	&	{\hr}kompo-\tbr{r}ii	{\cgr}&	\cg{(neeʔ)}	&	\cg{(pusi)}\\
Manggarai	&	{\hr}\tbr{d}ara	&	{\hr}\tbr{d}iʔa	&	{\hr}\tbr{d}eləm	&	\cg{(ʧəboŋ)}	&	{\hr}\tbr{s}iri	{\cgr}&	{\hr}\tbr{s}aki	{\cgr}&	{\hr}mo\tbr{s}e{\cgr}\\
Kepo{\Q}	&	{\hr}\tbr{d}ara	&	\cg{(wəru)}	&		&		&	{\hr}\tbr{d}iri	&	{\hr}\tbr{r}aki	{\cgr}&	{\hr}mo\tbr{r}e\\
Waerana	&	{\hr}\tbr{d}araʔ	&	\cg{(ŋgoʔi)}	&	\cg{(ləmaŋ)}	&	{\hr}\tbr{d}ioʔ	&	{\hr}\tbr{d}iri	&	{\hr}\tbr{d}aki	&	{\hr}mo\tbr{r}eʔ\\
Rembong	&	{\hr}\tbr{d}araʔ	&	{\hr}\tbr{d}iʔa	&	{\hr}\tbr{d}oloŋ	&	{\hr}\tbr{d}ioʔ	&	{\hr}\tbr{d}iri\su{†}	&	{\hr}\tbr{z}aki	{\cgr}&	{\hr}mo\tbr{r}eʔ\su{‡}\\
		\lspbottomrule
	\end{tabular}
		\begin{tablenotes}\tnsize
			\item[†]	{The Terong-Mawong variety of Rembong has \it{siri} ‘house post’.}
			\item[‡]	{The Wangka variety of Rembong has \it{muzit} ‘live’
								and Wue has \it{moseʔ} ‘live’.}
		\end{tablenotes}
	\end{threeparttable}
\end{adjustbox}
\end{table}

\begin{table}[p]\centering\tcp{\tsc{West Flores} *-aw}{03-32}
\begin{adjustbox}{width=\textwidth}
	\begin{threeparttable}\st{2pt}
	\begin{tabular}{llllll|ll}\lsptoprule
*gloss	&	sun, day	&	go	&	mouse	&	upper	&	Tinea flava	&	steal	&	\it{Grewia sp.}\\
PMP	&	*qaləj\tbr{aw}	&	*lak\tbr{aw}	&	*balab\tbr{aw}	&	*bab\tbr{aw}	&	*pan\tbr{aw}	&	*tak\tbr{aw}	&	*qanil\tbr{aw}\\ \midrule
Komodo	&	\hp{*qalə}r\tbr{oo}	&	{\hr}lah\tbr{o}	&	{\hr}belaw\tbr{o}	&	{\hr}wew\tbr{o}	&		&	{\hr}kənah\tbr{o}	&	{\hqa}nil\tbr{o}, nil\tbr{a}\su{†}\\
Manggarai	&	{\hqa}ləs\tbr{o}	&	{\hr}lak\tbr{o}	&	\hp{*ba}laʋ\tbr{o}	&	{\hr}wew\tbr{o}	&	{\hr}pan\tbr{o}	&	{\hr}tak\tbr{o}	&	{\hqa}nil\tbr{a}\\
Kepo{\Q}	&	{\hqa}ləz\tbr{o}ʔ	&	{\hr}lak\tbr{o}	&	\cg{(teʔus)}	&		&		&	{\hr}tak\tbr{o}	&	\\
Waerana	&	{\hqa}ləz\tbr{o}ʔ	&	{\hr}lak\tbr{o}ʔ	&	\cg{(teʔus)}	&	{\hr}waw\tbr{o}ʔ	&	{\hr}pan\tbr{o}	&	{\hr}tak\tbr{ar}	&	{\hqa}nil\tbr{ar}\\
Rembong	&	{\hqa}ləz\tbr{o}ʔ	&	{\hr}lak\tbr{o}ʔ	&	\cg{(endoŋ)}	&		&	{\hr}pan\tbr{o}	&	{\hr}tak\tbr{a}t	&	{\hqa}nil\tbr{a}\\
		\lspbottomrule
	\end{tabular}
		\begin{tablenotes}\tnsize
			\item[†]	{
								\citet[236]{Verheijen1987} gives Komodo \it{nila} as ‘\it{Grewia
								eriocarpa}’ (a kind of plant), while \citet[58]{Verheijen1984}
								gives Komodo \it{nilo} as ‘\it{Grewia
								spp.}’.}
		\end{tablenotes}
	\end{threeparttable}
\end{adjustbox}
\end{table}

\begin{table}[p]\centering\tcp{\tsc{West Flores} *-ay}{03-33}
%\begin{adjustbox}{width=\textwidth}
	\begin{threeparttable}\st{2pt}
	\begin{tabular}{lllllll}\lsptoprule
*gloss	&	die, dead	&	male	&	female	&	Job’s tears	&	liver	&	\it{Antidesma sp.}\\
PMP	&	*mat\tbr{ay}	&	*ma-Ruqan\tbr{ay}	&	*bin\tbr{ahi}	&	*zəl\tbr{ay}	&	*qat\tbr{ay}	&	*buRn\tbr{ay}\\ \midrule
Komodo	&	{\hr}mat\tbr{e}	&	{\hr}mon\tbr{e}	&	{\hr}win\tbr{e}	&	\cg{(gandoŋ)}	&	{\hq}at\tbr{e}	&	\\ 		\hhline{~----~~}
Manggarai	&	{\hr}mat\tbr{a}	&	{\hr}ron\tbr{a}	&	{\hr}ʋin\tbr{a}	&	\mc{1}{l|}{{\hr}sel\tbr{a}}	&	{\hq}at\tbr{i}	&	{\hr}ʋun\tbr{e}\\
Kepo{\Q}	&	{\hr}mat\tbr{a}	&	{\hr}ron\tbr{a}	&	{\hr}win\tbr{a}	&	\mc{1}{l|}{‹sel\tbr{a}›}	&	{\hq}at\tbr{i}	&	\\	\hhline{~~~~~-~}
Waerana	&	{\hr}mat\tbr{ar}	&	{\hr}ran\tbr{ar}	&	{\hr}win\tbr{ar}	&	{\hr}əl\tbr{ar}, əl\tbr{as}\su{†}	&	\mc{1}{l|}{{\hq}at\tbr{ar}}	&	{\hr}kazu wun\tbr{e}\\ \hhline{~~~~~-~}
Rembong	&	{\hr}mat\tbr{a}	&	{\hr}ran\tbr{a}	&	{\hr}win\tbr{a}	&	\mc{1}{l|}{\cg{(keʔo)}}	&	{\hq}at\tbr{i}	&	\\
		\lspbottomrule
	\end{tabular}
		\begin{tablenotes}\tnsize
			\item[†]	{
								\citet[14]{Verheijen1984} only gives Waerana \it{əlar} while
								\citet[11]{VerheijenDadu1994} give both \it{əlar} and \it{əlas}.
								The form \it{əlas} also occurs in Rajong \citep[14]{Verheijen1984}.}
		\end{tablenotes}
	\end{threeparttable}
%\end{adjustbox}
\end{table}
